\section*{UI Extension API}
\addcontentsline{toc}{section}{UI Extension API}

Symphony's UI Extension API enables out-of-process extensions to create rich user interfaces through a Virtual DOM bridge system. This approach maintains process isolation while providing full access to modern web UI capabilities.

\subsection*{Virtual DOM Bridge}

Extensions describe their UI using Virtual DOM structures that React efficiently renders:

\begin{lstlisting}[language=TypeScript,caption={Virtual DOM Interface}]
interface VirtualNode {
  type: string;
  props: Record<string, any>;
  children: VirtualNode[];
  key?: string;
  ref?: string;
}

interface UIExtensionAPI {
  // Render virtual DOM tree
  render(vdom: VirtualNode): Promise<void>;
  
  // Update specific nodes
  update(nodeId: string, props: Record<string, any>): Promise<void>;
  
  // Handle events from UI
  addEventListener(event: string, handler: EventHandler): void;
  
  // Remove event listeners
  removeEventListener(event: string, handler: EventHandler): void;
}
\end{lstlisting}

\subsection*{Component Registration}

Extensions register UI components with the Symphony component system:

\begin{lstlisting}[language=TypeScript,caption={Component Registration}]
interface ComponentRegistration {
  name: string;
  category: ComponentCategory;
  props: PropDefinition[];
  events: EventDefinition[];
  slots: SlotDefinition[];
  styling: StylingOptions;
}

enum ComponentCategory {
  Panel = "panel",
  Editor = "editor",
  Sidebar = "sidebar",
  Modal = "modal",
  Toolbar = "toolbar",
  StatusBar = "status_bar"
}

interface PropDefinition {
  name: string;
  type: PropType;
  required: boolean;
  default?: any;
  description: string;
}

// Register component with Symphony
const registration: ComponentRegistration = {
  name: "MyCustomEditor",
  category: ComponentCategory.Editor,
  props: [
    {
      name: "content",
      type: PropType.String,
      required: true,
      description: "Editor content"
    },
    {
      name: "language",
      type: PropType.String,
      required: false,
      default: "plaintext",
      description: "Syntax highlighting language"
    }
  ],
  events: [
    {
      name: "contentChanged",
      payload: { content: PropType.String },
      description: "Fired when content changes"
    }
  ],
  slots: [],
  styling: {
    themes: ["light", "dark"],
    customizable: true
  }
};
\end{lstlisting}

\subsection*{Event Handling}

UI extensions handle user interactions through a structured event system:

\begin{lstlisting}[language=TypeScript,caption={Event Handling System}]
interface UIEvent {
  type: string;
  target: string;
  data: Record<string, any>;
  timestamp: number;
  preventDefault?: boolean;
  stopPropagation?: boolean;
}

interface EventHandler {
  (event: UIEvent): Promise<void> | void;
}

// Example event handlers
class MyExtensionUI {
  private api: UIExtensionAPI;
  
  constructor(api: UIExtensionAPI) {
    this.api = api;
    this.setupEventHandlers();
  }
  
  private setupEventHandlers() {
    // Handle button clicks
    this.api.addEventListener('click', async (event) => {
      if (event.target === 'save-button') {
        await this.handleSave(event.data);
      }
    });
    
    // Handle text input
    this.api.addEventListener('input', async (event) => {
      if (event.target === 'editor-content') {
        await this.handleContentChange(event.data.value);
      }
    });
    
    // Handle keyboard shortcuts
    this.api.addEventListener('keydown', async (event) => {
      if (event.data.ctrlKey && event.data.key === 's') {
        event.preventDefault = true;
        await this.handleSave();
      }
    });
  }
  
  private async handleSave(data?: any) {
    // Implementation
  }
  
  private async handleContentChange(content: string) {
    // Implementation
  }
}
\end{lstlisting}

\subsection*{Layout System}

Extensions can create complex layouts using Symphony's layout primitives:

\begin{lstlisting}[language=TypeScript,caption={Layout System}]
interface LayoutNode extends VirtualNode {
  layout: LayoutProperties;
}

interface LayoutProperties {
  display: 'flex' | 'grid' | 'block' | 'inline';
  direction?: 'row' | 'column';
  justify?: 'start' | 'center' | 'end' | 'space-between' | 'space-around';
  align?: 'start' | 'center' | 'end' | 'stretch';
  gap?: number;
  padding?: Spacing;
  margin?: Spacing;
  width?: Size;
  height?: Size;
}

interface Spacing {
  top?: number;
  right?: number;
  bottom?: number;
  left?: number;
}

type Size = number | 'auto' | 'fit-content' | `${number}%` | `${number}px`;

// Example layout
const editorLayout: LayoutNode = {
  type: 'div',
  props: { className: 'editor-container' },
  layout: {
    display: 'flex',
    direction: 'column',
    height: '100%'
  },
  children: [
    {
      type: 'div',
      props: { className: 'toolbar' },
      layout: {
        display: 'flex',
        direction: 'row',
        justify: 'space-between',
        padding: { top: 8, bottom: 8, left: 16, right: 16 }
      },
      children: [/* toolbar items */]
    },
    {
      type: 'div',
      props: { className: 'editor-content' },
      layout: {
        display: 'block',
        height: 'auto'
      },
      children: [/* editor content */]
    }
  ]
};
\end{lstlisting}

\subsection*{Theming and Styling}

Extensions integrate with Symphony's theming system:

\begin{lstlisting}[language=TypeScript,caption={Theming System}]
interface ThemeAPI {
  // Get current theme
  getCurrentTheme(): Promise<Theme>;
  
  // Subscribe to theme changes
  onThemeChanged(callback: (theme: Theme) => void): void;
  
  // Get theme variables
  getThemeVariable(name: string): Promise<string>;
  
  // Apply custom styles
  applyStyles(styles: StyleSheet): Promise<void>;
}

interface Theme {
  name: string;
  type: 'light' | 'dark';
  colors: ThemeColors;
  typography: Typography;
  spacing: SpacingScale;
  shadows: ShadowScale;
}

interface ThemeColors {
  primary: string;
  secondary: string;
  accent: string;
  background: string;
  surface: string;
  text: string;
  textSecondary: string;
  border: string;
  error: string;
  warning: string;
  success: string;
  info: string;
}

// Example theme usage
class ThemedComponent {
  private theme: Theme;
  
  async render() {
    this.theme = await themeAPI.getCurrentTheme();
    
    const styles = {
      container: {
        backgroundColor: this.theme.colors.surface,
        color: this.theme.colors.text,
        border: `1px solid ${this.theme.colors.border}`,
        borderRadius: '4px',
        padding: this.theme.spacing.medium
      },
      button: {
        backgroundColor: this.theme.colors.primary,
        color: this.theme.colors.background,
        border: 'none',
        borderRadius: '4px',
        padding: `${this.theme.spacing.small} ${this.theme.spacing.medium}`
      }
    };
    
    return {
      type: 'div',
      props: { style: styles.container },
      children: [
        {
          type: 'button',
          props: { 
            style: styles.button,
            onClick: 'handleClick'
          },
          children: ['Click me']
        }
      ]
    };
  }
}
\end{lstlisting}

\subsection*{State Management}

Extensions can manage UI state through Symphony's state management system:

\begin{lstlisting}[language=TypeScript,caption={State Management}]
interface StateManager {
  // Get current state
  getState<T>(key: string): Promise<T | undefined>;
  
  // Set state
  setState<T>(key: string, value: T): Promise<void>;
  
  // Subscribe to state changes
  subscribe<T>(key: string, callback: (value: T) => void): () => void;
  
  // Batch state updates
  batch(updates: Record<string, any>): Promise<void>;
}

// Example state management
class StatefulComponent {
  private stateManager: StateManager;
  private unsubscribe: (() => void)[] = [];
  
  constructor(stateManager: StateManager) {
    this.stateManager = stateManager;
    this.setupStateSubscriptions();
  }
  
  private setupStateSubscriptions() {
    // Subscribe to editor content changes
    const unsubContent = this.stateManager.subscribe<string>(
      'editor.content',
      (content) => this.handleContentChange(content)
    );
    this.unsubscribe.push(unsubContent);
    
    // Subscribe to selection changes
    const unsubSelection = this.stateManager.subscribe<Selection>(
      'editor.selection',
      (selection) => this.handleSelectionChange(selection)
    );
    this.unsubscribe.push(unsubSelection);
  }
  
  private async handleContentChange(content: string) {
    // Update UI based on content change
    await this.api.update('editor-content', { value: content });
  }
  
  private async handleSelectionChange(selection: Selection) {
    // Update UI based on selection change
    await this.api.update('selection-info', { 
      text: `Line ${selection.line}, Column ${selection.column}` 
    });
  }
  
  public cleanup() {
    this.unsubscribe.forEach(fn => fn());
  }
}
\end{lstlisting}

\subsection*{Performance Optimization}

The UI Extension API includes performance optimization features:

\begin{lstlisting}[language=TypeScript,caption={Performance Optimization}]
interface PerformanceAPI {
  // Batch DOM updates
  batchUpdates(updates: () => void): Promise<void>;
  
  // Virtualize large lists
  createVirtualList(options: VirtualListOptions): VirtualList;
  
  // Memoize expensive computations
  memoize<T>(fn: (...args: any[]) => T, deps: any[]): T;
  
  // Debounce frequent updates
  debounce<T extends (...args: any[]) => any>(fn: T, delay: number): T;
  
  // Throttle high-frequency events
  throttle<T extends (...args: any[]) => any>(fn: T, interval: number): T;
}

interface VirtualListOptions {
  itemHeight: number;
  containerHeight: number;
  items: any[];
  renderItem: (item: any, index: number) => VirtualNode;
  overscan?: number;
}

// Example performance optimization
class OptimizedList {
  private performanceAPI: PerformanceAPI;
  private debouncedSearch: (query: string) => void;
  
  constructor(performanceAPI: PerformanceAPI) {
    this.performanceAPI = performanceAPI;
    
    // Debounce search to avoid excessive API calls
    this.debouncedSearch = this.performanceAPI.debounce(
      this.performSearch.bind(this),
      300
    );
  }
  
  private async performSearch(query: string) {
    // Expensive search operation
    const results = await this.searchAPI.search(query);
    
    // Batch UI updates for better performance
    await this.performanceAPI.batchUpdates(() => {
      this.updateSearchResults(results);
      this.updateResultCount(results.length);
      this.updateLoadingState(false);
    });
  }
  
  private createVirtualizedList(items: any[]) {
    return this.performanceAPI.createVirtualList({
      itemHeight: 40,
      containerHeight: 400,
      items: items,
      renderItem: (item, index) => ({
        type: 'div',
        props: { 
          key: item.id,
          className: 'list-item'
        },
        children: [item.name]
      }),
      overscan: 5
    });
  }
}
\end{lstlisting}

This UI Extension API enables rich, performant user interfaces while maintaining the security and isolation benefits of Symphony's out-of-process extension architecture.
